\documentclass[11pt,a4paper]{article}
\usepackage[T1]{fontenc}
\usepackage{lmodern,microtype}
\usepackage[margin=27mm]{geometry}
\usepackage{amsmath,amssymb,amsthm,mathtools,booktabs,enumitem,xcolor}
\usepackage[colorlinks=true,linkcolor=blue!45!black,citecolor=blue!45!black,urlcolor=blue!45!black]{hyperref}
\newcommand{\N}{\mathbb N}\newcommand{\Z}{\mathbb Z}\newcommand{\E}{\mathbb E}\newcommand{\Prob}{\mathbb P}
\newcommand{\ind}{\mathbf1}\newcommand{\spn}{\operatorname{span}}\newcommand{\MPF}{\operatorname{MPF}}\newcommand{\mpf}{\operatorname{mpf}}\newcommand{\Inj}{\operatorname{Inj}}
\newcommand{\pos}[1]{\left(#1\right)_+}\newcommand{\clipbin}[2]{\binom{#1}{#2}_{\!+}}
\newtheorem{theorem}{Theorem}[section]\newtheorem{lemma}[theorem]{Lemma}\newtheorem{corollary}[theorem]{Corollary}\newtheorem{proposition}[theorem]{Proposition}\theoremstyle{remark}\newtheorem{remark}[theorem]{Remark}\newtheorem{example}[theorem]{Example}
\setlength{\parskip}{3pt}\setlength{\parindent}{1.5em}\setlist{nosep}\allowdisplaybreaks[2]\emergencystretch=2em
\hypersetup{pdftitle={Metered Parking Functions and Prescribed Lucky Cars},pdfauthor={Peilin Liu},pdfsubject={Combinatorics}}
\title{Metered Parking Functions and Prescribed Lucky Cars}
\author{Peilin Liu\\\small Shandong University\\\small\texttt{202617552@mail.sdu.edu.cn}}
\date{September 29, 2026}
\begin{document}\maketitle
\begin{abstract}
We enumerate metered parking functions by their lucky cars, allowing arbitrary predetermined departure schedules. The fibers of the parking outcome map yield a uniform affine formula valid for every street size and a positive expansion in packed outcomes. For a prescribed lucky set, we prove eventual polynomiality, determine the degree and leading coefficient, and obtain a graphical bound on component spans. For a common meter this bound is attained and depends explicitly on the positions of the unlucky cars. Summing the leading coefficients gives the chromatic polynomial of the overlap graph; the positive expansion gives nonnegative Newton coefficients. We also give an exact algorithm that is polynomial at fixed occupancy width. For meter one, a reduced rational generating function and a positive recurrence count every number of cars, including more cars than sites, while an interval-partition formula determines the exact polynomial onset for every prescribed lucky set. Finally, a finite boundary comparison identifies the highest coefficients of the total counting polynomial, including the next coefficient that detects chronology.
\end{abstract}
\noindent\textbf{Keywords:} metered parking functions; lucky cars; outcome fibers; affine enumeration; interval graphs.\\
\textbf{2020 Mathematics Subject Classification:} Primary 05A15; Secondary 05A05, 05C15.

\section{Introduction}
Daugherty, Harris, Klein, and McClinton introduced metered parking functions and enumerated several special regimes \cite{DH}. Cars arrive in a fixed order, search to the right from their preferred sites, and later depart. Departures allow more cars than sites to succeed, but they do not undo the displacements caused by earlier arrivals. Thus ordinary parking formulas do not directly extend to this setting.

The central question is to count successful preference words while retaining the arrival-order information. We refine this question by prescribing the set of lucky cars, namely those that park at their preferred sites. Our starting point is an outcome word: once all parking positions are fixed, the preferences of each car form an interval determined by the occupied block immediately below its outcome. This fiber description yields two complementary finite formulas. One records integer difference constraints and applies inclusion--exclusion; the other packs the used sites into consecutive runs and has nonnegative coefficients.

The principal results are as follows.
\begin{enumerate}
\item For any predetermined departure schedule, Theorem~\ref{thm:affine} gives an exact affine enumeration, with lucky cars marked individually. Theorem~\ref{thm:fixed} proves that a prescribed lucky set $K$ has a stable polynomial of degree $|K|$ and leading coefficient $\prod_{i\notin K}r_i$, where $r_i$ is the number of cars active before arrival $i$.
\item For a common meter $t$, the maximum fixed-set component span is exactly the number of adjacent time edges covered by the windows ending at unlucky arrivals. This refines the bound $\min\{m-1,t(m-|K|)\}$ and gives an explicit range of polynomial agreement. The chromatic identity in Corollary~\ref{cor:leading} records the leading data of all lucky layers.
\item Theorem~\ref{thm:packed} gives a positive packed-outcome expansion and nonnegative Newton coefficients. Proposition~\ref{prop:algorithm} compresses the affine formula into an exact algorithm that is polynomial for every fixed occupancy width.
\item For meter one, Theorem~\ref{thm:one} gives a reduced rational generating function, positive recurrence, and finite multinomial formula valid for every $m,n$. The prescribed-run formula in Section~\ref{sec:one} also determines the smallest polynomial onset for every lucky set.
\end{enumerate}
These results address the general-meter enumeration and lucky-layer degree questions, as well as the meter-one enumeration for $m>n$, posed in the published version of \cite{DH}. The general affine formula is structural and may contain exponentially many terms; the meter-one formulas and fixed-width algorithm provide separate effective specializations. A short boundary analysis in Section~\ref{sec:boundary} relates the highest coefficients to the overlap graph and explains where chronology first enters.

Affine lattice counting and gain-graph inclusion--exclusion are established methods \cite{FZ,BC}. Here the parking-specific content is the weighted constraint alphabet obtained from the fibers, the attained common-meter span, and the leading coefficients for prescribed lucky sets. Fixed-lucky-set parking without departures was studied by Harris and Martinez \cite{HM}, and capacity-based refinements by Ferreri, Harris, Martinez, and Swartz \cite{FHMS}. The present departure model requires the active set to change with arrival time; a single outcome may reuse a site after its occupant leaves.

\section{Outcomes and their preference fibers}\label{sec:outcomes}

Write $[m]=\{1,\ldots,m\}$ and $[0]=\varnothing$.
Fix $m\ge 0$, $n\ge 1$, and integers $d_j\ge j$ for $j\in[m]$.
Car $i$, with preference $a_i\in[n]$, parks at the first vacant spot
weakly to the right of $a_i$.  If no such spot exists, it fails.  Immediately
after the search at time $i$, all cars whose departure index is $i$ leave.
Departures after time $m$ have no effect, so the indices $d_j$ may equivalently
be replaced by $\min(d_j,m)$.
Before the $i$th search, the active earlier cars are
\[
 J_i=\{j<i:d_j\ge i\},\qquad r_i=|J_i|,\qquad
 R=\max_{1\le i\le m}r_i.
\]
Set $R=0$ when $m=0$.  A common meter $t\ge0$ means $d_j=j+t$, and then
$J_i=\{\max(1,i-t),\ldots,i-1\}$ and $r_i=\min(t,i-1)$.
A successful word has an outcome $p=(p_1,\ldots,p_m)$, where $p_i$ is the
spot chosen by car $i$.  Car $i$ is lucky if $a_i=p_i$.
Let $\MPF_{d,n}$ be the set of successful preference words for this schedule. Write
\[
 L_{d,n}(u)=\sum_{a\in\MPF_{d,n}}u^{|\operatorname{Lucky}(a)|},
 \qquad
 N_{d,K}(n)=\#\{a\in\MPF_{d,n}:\operatorname{Lucky}(a)=K\}.
\]
For a common meter write $\MPF_{m,n}(t)=\MPF_{d,n}$, $L_{m,n,t}$, $\mpf_{m,n}(t)=L_{m,n,t}(1)$, and $N_{m,t,K}(n)=N_{d,K}(n)$. Empty-word counts equal one.  The overlap graph $G_d$ has vertex set $[m]$
and edges $\{j,i\}$ for $j\in J_i$.  It is the intersection graph of the
intervals $[j,d_j]$.  In particular, $J_i\cup\{i\}$ is a clique.

\begin{lemma}[Outcome characterization]\label{lem:outcome}
An outcome word $p\in[n]^m$ is feasible if and only if
\[
 p_i\ne p_j\qquad(i\in[m],\ j\in J_i).
\]
For $m\ge1$, successful words exist if and only if $n\ge R+1$.
For a common meter this condition is $n\ge\min(m,t+1)$.
\end{lemma}
\begin{proof}
If a word succeeds, the active cars occupy distinct spots, so the displayed
inequalities are necessary.  Conversely, suppose they hold and choose
$a_i=p_i$ at every time.  Inductively, when car $i$ arrives, spot $p_i$ is
unoccupied by every member of $J_i$; the car parks there without moving.
Thus every proper $[n]$-coloring of $G_d$ is realized as an outcome.

Choose $i$ with $r_i=R$.  The clique $J_i\cup\{i\}$ requires $R+1$ distinct
spots, proving necessity of the bound.  For sufficiency, color the vertices in
chronological order.  At vertex $i$, only the $r_i\le R$ earlier neighbors
forbid colors, so $R+1$ colors suffice.  For common meters,
$R=\min(t,m-1)$, giving the stated specialization.
\end{proof}

For a feasible outcome, let $b_i(p)$ be the largest integer $h\ge0$ for
which $p_i-1,\ldots,p_i-h$ are all occupied immediately before car $i$
arrives.  Only positive spots are included, so $b_i(p)\le p_i-1$.
The preferences producing $p_i$ are exactly
\[
 p_i,p_i-1,\ldots,p_i-b_i(p).
\]
Indeed, all the intermediate spots are occupied for these preferences,
whereas starting further left encounters a vacant spot before $p_i$.
The lucky choice is the unique preference $p_i$, and the other $b_i(p)$
choices are unlucky.  Since the departure schedule is fixed, specifying all
the outcomes fixes the active set at each arrival independently of the
preference choices.  Consequently, the lucky polynomial of the fiber over
$p$ is $\prod_i(u+b_i(p))$.

Let $\Inj([h],J_i)$ denote the injections from $[h]$ into $J_i$, and put
\[
 W_i(p;u)=u+\sum_{h=1}^{r_i}\sum_{\phi\in\Inj([h],J_i)}
             \prod_{q=1}^h\ind\{p_i-p_{\phi(q)}=q\}.
\]
\begin{lemma}[Marked fiber identity]\label{lem:fiber}
For every $m\ge0$ and $n\ge1$,
\[
 L_{d,n}(u)=\sum_{p\in[n]^m}
    \left(\prod_i\prod_{j\in J_i}(1-\ind\{p_i=p_j\})\right)
    \prod_iW_i(p;u).
\]
Replacing the first term of $W_i$ by a separate variable $u_i$, the coefficient
of $\prod_{i\in K}u_i$ in this identity equals $N_{d,K}(n)$.
\end{lemma}
\begin{proof}
The first product is the feasibility indicator from
Lemma~\ref{lem:outcome}.  For a feasible outcome and $1\le h\le b_i(p)$,
there is exactly one injection $\phi$ with the displayed equations: the
distinct active outcomes uniquely identify the cars at the spots
$p_i-1,\ldots,p_i-h$.  If $h>b_i(p)$, there is no such injection.  Hence
$W_i(p;u)=u+b_i(p)$, and the preceding fiber calculation proves the identity.
With separate variables, choosing $u_i$ specifies that car $i$ is lucky;
choosing one of the other terms specifies one of its unlucky preferences.
Each variable has degree at most one, proving the coefficient statement.
\end{proof}

\section{Affine enumeration and prescribed lucky sets}\label{sec:exact}
\subsection{The affine representation}

At time $i$, a displacement decoration is either empty or a pair
$(\ell_i,\phi_i)$ with $1\le\ell_i\le r_i$ and
$\phi_i\in\Inj([\ell_i],J_i)$.  A nonempty decoration imposes
\[
 p_i-p_{\phi_i(q)}=q\qquad(1\le q\le\ell_i).
\]
Independently choose $B_i\subseteq J_i$.  It imposes the zero-difference
constraints $p_i-p_j=0$ for $j\in B_i$ and contributes the sign
$(-1)^{|B_i|}$.  Let $M_d$ be the set of all indexed collections of these
local choices.  For $\mu\in M_d$, write
\[
 K(\mu)=\{i:\text{the displacement decoration at $i$ is empty}\},
 \qquad \epsilon(\mu)=(-1)^{\sum_i|B_i|}.
\]
The constraint graph of $\mu$ has vertex set $[m]$ and one edge for every
selected equation.  Multiple selected equations are retained: contradictory
equations make the system inconsistent.  In a consistent connected component
$C$, there are integer potentials $h_v$ such that all its solutions have the
form $p_v=z_C+h_v$.  These potentials are unique up to a common additive
constant.  Define
\[
 \spn_\mu(C)=\max_{v\in C}h_v-\min_{v\in C}h_v,
 \qquad
 \Phi_n(\mu)=
 \begin{cases}
  \displaystyle\prod_C(n-\spn_\mu(C))_+,&\mu\text{ consistent},\\
  0,&\mu\text{ inconsistent}.
 \end{cases}
\]
Isolated vertices are components of span zero; empty products equal one.

\begin{theorem}[Uniform affine enumeration]\label{thm:affine}
For every fixed departure schedule and all $n\ge1$,
\[
 L_{d,n}(u)=\sum_{\mu\in M_d}\epsilon(\mu)u^{|K(\mu)|}\Phi_n(\mu),
 \qquad
 N_{d,K}(n)=\sum_{\substack{\mu\in M_d\\K(\mu)=K}}
                    \epsilon(\mu)\Phi_n(\mu).
\]
The number of formal indexed terms is $\prod_i c_{r_i}$, where
\[
 c_s=2^s\sum_{h=0}^s\frac{s!}{(s-h)!}.
\]
In particular $c_0,c_1,c_2,c_3=1,4,20,128$.
\end{theorem}
\begin{proof}
Expand the products in Lemma~\ref{lem:fiber}.  Choosing $u$ or one summand
from $W_i$ gives exactly the displacement decoration; expanding the
feasibility product gives exactly $B_i$ and its sign.  A term is zero unless
all its selected equations hold.

For a consistent component $C$, the restrictions $1\le z_C+h_v\le n$ for
all $v\in C$ are equivalent to
\[
 1-\min_{v\in C}h_v\le z_C\le n-\max_{v\in C}h_v.
\]
The number of integer choices is $(n-\spn_\mu(C))_+$.  Different components
have no equations between them and can be translated independently.
Thus the number of outcome assignments satisfying a selected system is
$\Phi_n(\mu)$.  This gives the first identity, and the separate-variable
version of Lemma~\ref{lem:fiber} gives the fixed-set identity.

At a time with $r_i=s$, there are $s!/(s-h)!$ injections of length $h$;
the summand $h=0$ represents the unique empty decoration.  There are
$2^s$ independent equality choices.  Multiplying the local counts proves
the formula for the number of indexed terms.
\end{proof}

\subsection{Span bounds and fixed-set leading coefficients}

Let $Z=\{i:r_i=0\}$.  Given a proposed lucky set $K\subseteq[m]$, put
$U=[m]\setminus K$, and let $H_{d,U}$ be the graph on all $m$ vertices
obtained by taking the union of the cliques on $\{i\}\cup J_i$ for $i\in U$.
Isolated vertices are counted when writing $c(H_{d,U})$.  Its graphical
rank is
\[
 D_{d,K}=m-c(H_{d,U}).
\]
For $m\ge1$,
\[
 0\le D_{d,K}\le\min\left(m-1,\sum_{i\in U}r_i\right).
\]
The second bound follows because each added clique has rank $r_i$, and
adding its edges increases the rank of the union by at most $r_i$.
Define the integer polynomial
\[
 Q_{d,K}(x)=\sum_{\substack{\mu\in M_d:\ K(\mu)=K\\
                            \mu\ {\rm consistent}}}
                \epsilon(\mu)\prod_C(x-\spn_\mu(C)).
\]
For a common meter write $Q_{m,t,K}=Q_{d,K}$.

\begin{theorem}[Prescribed lucky-set polynomials]\label{thm:fixed}
If $Z\not\subseteq K$, then $N_{d,K}=Q_{d,K}=0$.
If $Z\subseteq K$, then
\[
 N_{d,K}(n)=Q_{d,K}(n)\quad(n\ge\max(1,D_{d,K})),
 \qquad
 \deg Q_{d,K}=|K|,
 \qquad
 [x^{|K|}]Q_{d,K}(x)=\prod_{i\notin K}r_i.
\]
For $m=0$, the assertion is interpreted as $N_{d,\varnothing}=Q_{d,\varnothing}=1$.
\end{theorem}
\begin{proof}
If $i\in Z$, no nonempty displacement decoration is available at time $i$;
equivalently, a car arriving to an empty street must be lucky.  Thus the
fixed-set sum is empty whenever $Z\not\subseteq K$.

For any remaining consistent motif, replace each nonempty displacement star
by the equations
\[
 p_i-p_{\phi_i(1)}=1,\qquad
 p_{\phi_i(q)}-p_{\phi_i(q+1)}=1\quad(1\le q<\ell_i).
\]
This unit-chain replacement has exactly the same solutions and the same
connected components as the original star.  All its nonzero-difference edges
belong to $H_{d,U}$, because both their endpoints lie in
$\{i\}\cup J_i$ for the relevant $i\in U$.

In one connected constraint component $C$, take a simple path from a
vertex of minimum potential to a vertex of maximum potential.  Each
nonzero-difference edge changes the potential by absolute value one;
all other edges change it by zero.  The nonzero edges of this simple path
form a forest in $H_{d,U}$, so there are at most $D_{d,K}$ of them.
This remains true even when zero-difference edges on the path connect
different components of $H_{d,U}$.  Consequently
\[
 \spn_\mu(C)\le D_{d,K}.
\]
The same path has at most $|C|-1$ edges, giving the useful independent
bound $\spn_\mu(C)\le|C|-1$.
For $n\ge D_{d,K}$, every factor $(n-\spn_\mu(C))_+$ equals
$n-\spn_\mu(C)$, including the boundary case of a zero factor.
Theorem~\ref{thm:affine} therefore proves the polynomial identity in the
stated range.

Every connected constraint component contains an index in $K$.
Indeed, if its least index $i$ were in $U$, its nonempty displacement
would connect it to an earlier index $\phi_i(1)\in J_i$, a contradiction.
Thus a fixed-$K$ motif has at most $|K|$ connected components, and its
polynomial contribution has degree at most $|K|$.

It remains to identify the leading coefficient without relying on
cancellation in the signed sum.  Suppose $R>0$ and $m\ge1$, and write
$k=|K|$ and $C=mR$.  An instruction consists of a root value
$\beta_i\in[n]$ for each $i\in K$ and a parent
$j_i\in J_i$ for each $i\in U$.  Its preferences are constructed in time
order by
\[
 a_i=\beta_i\quad(i\in K),\qquad a_i=p_{j_i}\quad(i\in U).
\]
Call the root tuple generic if
\[
 1\le\beta_i\le n-C\quad(i\in K),\qquad
 |\beta_i-\beta_j|>2C\quad(i\ne j\text{ in }K).
\]
For fixed $m$ and schedule, the number $G(n)$ of generic tuples satisfies
$G(n)=n^k+O(n^{k-1})$.  Explicitly, boundary violations account for at
most $kC n^{k-1}$ tuples, and close pairs account for at most
$\binom{k}{2}(4C+1)n^{k-1}$ tuples, for all sufficiently large $n$.

Each generic tuple and every parent selection produce a successful word
with lucky set exactly $K$.  To see this, first run the construction on an
unbounded street.  Following parent indices strictly decreases time and
ends at a root in $K$.  Inductively, an outcome at parent depth $s$ from
root $\beta$ lies in $[\beta,\beta+sR]$: an arrival sees at most $R$
occupied spots and therefore moves at most $R$ spots.  At a root time
$i$, all earlier outcomes lie in the intervals of earlier roots; these
intervals are disjoint from $\beta_i$, so car $i$ is lucky.  At an
unlucky time, the parent is still active, so the preferred spot is
occupied and the arriving car moves strictly right.  All outcomes lie
in the disjoint root intervals $[\beta,\beta+C]\subseteq[n]$.
Thus the construction succeeds on the finite street as well.

Conversely, a successful word with lucky set $K$ uniquely determines its
instructions: at a lucky time $i$, $\beta_i=a_i=p_i$; at an unlucky time
$i$, the preference $a_i$ is occupied by exactly one active car, which
uniquely determines $j_i$.  The instructions then reconstruct the word
chronologically.  In particular, the number of words with any root
tuple is at most the fixed number $M=\prod_{i\in U}r_i$ of parent
selections.  Nongeneric tuples therefore contribute at most
$M(n^k-G(n))=O(n^{k-1})$ words, whereas generic tuples contribute
exactly $MG(n)$ words.  It follows that
\[
 N_{d,K}(n)=\left(\prod_{i\in U}r_i\right)n^k+O(n^{k-1}).
\]
Since $r_i>0$ for $i\in U$ and the count eventually agrees with
$Q_{d,K}$, this identifies a positive leading coefficient and proves
the degree assertion.  If $R=0$, then $Z=[m]$, the only possible set
is $K=[m]$, and $N_{d,K}(n)=n^m$; its polynomial is $x^m$.
\end{proof}

\subsection{An attained position-sensitive span for common meters}

For a common meter $t\ge1$, fix $m\ge1$ and $K\subseteq[m]$ with $1\in K$.
Put $U=[m]\setminus K$ and
\[
 \mathcal E_t(U)=\bigcup_{i\in U}\{\max(2,i-t+1),\ldots,i\},
 \qquad D_t(U)=|\mathcal E_t(U)|.
\]
An index $j\in\mathcal E_t(U)$ represents the adjacent pair of arrival
times $j-1,j$.  If $U=\{u_1<\cdots<u_q\}$ is nonempty, the same quantity is
\[
 D_t(U)=\min(t,u_1-1)+\sum_{a=2}^{q}\min(t,u_a-u_{a-1});
 \qquad D_t(\varnothing)=0.
\]
This expression follows by adding the intervals in increasing order of
their right endpoints.

\begin{theorem}[Exact fixed-set span]\label{thm:fixedspan}
The largest component span among consistent motifs with lucky set $K$
is exactly $D_t(U)$.  Hence
\[
 N_{m,t,K}(n)=Q_{m,t,K}(n)\qquad(n\ge\max(1,D_t(U))).
\]
For $1\le k\le m$, the largest span over lucky sets of size $k$ is
\[
 \rho_{m,t,k}=\min(m-1,t(m-k)).
\]
\end{theorem}
\begin{proof}
The graph $H_{d,U}$ is the union of the cliques on
$\{\max(1,i-t),\ldots,i\}$ for $i\in U$.
Replacing each such clique by its adjacent path preserves connectivity.
The resulting graph is a subgraph of the path on $[m]$, with edge set
$\{\{j-1,j\}:j\in\mathcal E_t(U)\}$.  It is a forest, so
\[
 D_{d,K}=m-c(H_{d,U})=D_t(U).
\]
The preceding theorem bounds every component span by $D_t(U)$.

For attainment, set $h_1=0$ and
\[
 h_j-h_{j-1}=\ind\{j\in\mathcal E_t(U)\}\qquad(2\le j\le m).
\]
At each $i\in U$, select the full displacement decoration of length
$r_i=\min(t,i-1)$ with $\phi_i(q)=i-q$.
All adjacent increments in its time window belong to $\mathcal E_t(U)$,
so $h_i-h_{i-q}=q$ for $1\le q\le r_i$; all its equations are consistent
with the chosen potentials.  Select the empty displacement at every
time in $K$.  For each $j\ge2$ outside $\mathcal E_t(U)$, select the
zero-difference edge $\{j-1,j\}$, and select no other zero edges.
This is an allowed equality choice because $t\ge1$ makes $j-1$ active
at time $j$, and it is consistent because $h_j=h_{j-1}$.
After the unit-chain replacement, every adjacent pair is connected by
either a unit-difference edge or a zero-difference edge.  The motif
therefore has one component, and its span is
$h_m-h_1=D_t(U)$.  Notice that this is a formal consistent motif in
the inclusion--exclusion expansion; it need not itself be a feasible
parking outcome.

Finally let $q=m-k$.  Any union of $q$ windows has at most $tq$ edges
and at most $m-1$ edges.  If $q>0$, put $\rho=\min(m-1,tq)$.
Since $q\le\rho\le tq$, choose positive integers
$\ell_1,\ldots,\ell_q\le t$ summing to $\rho$ and set
$u_a=1+\ell_1+\cdots+\ell_a$.  These indices lie in
$\{2,\ldots,m\}$, and their windows form a connected union covering
exactly the edges with indices $2,\ldots,\rho+1$.
Thus $D_t(\{u_1,\ldots,u_q\})=\rho$.  The case $q=0$ has span zero.
\end{proof}

The fixed-set bound is optimal for removing positive parts term by term:
the attained motif has the single factor $(n-D_t(U))_+$.
This does not assert that $D_t(U)$ is the smallest onset for the signed
sum $N_{m,t,K}$.  Different motifs can cancel below this bound.
For instance, $m=5$, $t=2$, and $U=\{2,3\}$ give
$D_t(U)=2$, while $\sum_{i\in U}r_i=3$ and $t|U|=4$.
For general schedules the graphical-rank bound itself need not be attained:
if $d=(3,2,3)$ and $K=\{1\}$, then $D_{d,K}=2$, but every relevant
consistent motif has span at most one and $N_{d,K}(n)=n-1$.

\begin{corollary}[Lucky layers and chromatic leading data]\label{cor:leading}
Let $Q_{d,k}(x)=\sum_{|K|=k}Q_{d,K}(x)$, $Q_d=\sum_kQ_{d,k}$, and
$z_0=|Z|$. For a common meter write $Q_{m,t}=Q_d$.
Then $Q_{d,k}=0$ for $k<z_0$ and $\deg Q_{d,k}=k$ for
$z_0\le k\le m$.  These polynomials give the exact lucky layers and
total count whenever $n\ge\max(1,m-1)$, and
\[
 \sum_k[x^k]Q_{d,k}(x)q^k
       =\prod_{i=1}^m(q+r_i)=(-1)^m\chi_{G_d}(-q).
\]
The total polynomial is monic of degree $m$.  For a common meter, the
$k$th layer already agrees with its polynomial on
$n\ge\max(1,\rho_{m,t,k})$.
\end{corollary}
\begin{proof}
Sum Theorem~\ref{thm:fixed} over lucky sets.  All eligible leading
coefficients are positive, and their sum is the coefficient of $q^k$
in $\prod_i(q+r_i)$.  The factors with $r_i=0$ force the minimum
degree $z_0$.  Chronological coloring of $G_d$ gives
$\chi_{G_d}(x)=\prod_i(x-r_i)$: at time $i$, the $r_i$ earlier
neighbors form a clique and therefore forbid exactly $r_i$ distinct
colors.  The only degree-$m$ lucky layer is $K=[m]$, with leading
coefficient one.  The polynomial ranges follow from the span bounds
proved above.
\end{proof}


\subsection{Positive packing and a nonnegative Newton basis}

Throughout this subsection let $m\ge1$. For $1\le s\le m$, take a
surjective word $v\in[s]^m$ that is proper for $G_d$, and a set of cuts
$C\subseteq[s-1]$. A cut at $a$ separates labels $a$ and $a+1$.
Define $\beta_i(v,C)$ to be the largest $h\ge0$ for which, for every
$1\le q\le h$,
\[
 v_i-q\ge1,\qquad v_i-q\in\{v_j:j\in J_i\},\qquad v_i-q\notin C.
\]
Thus the block below $v_i$ stops at either an absent active label or a
cut. For $1\le b\le s$ set
\[
 C_{d;s,b}(u)=
 \sum_{\substack{v\in[s]^m\text{ surjective and proper}\\
                  C\subseteq[s-1],\ |C|=b-1}}
 \prod_{i=1}^m\bigl(u+\beta_i(v,C)\bigr).
\]
For integers $N$ and $b\ge1$, use the clipped convention
$\clipbin{N}{b}=\binom Nb$ if $N\ge b$, and zero otherwise.

\begin{theorem}[Positive packed enumeration]\label{thm:packed}
For every $n\ge1$,
\[
 L_{d,n}(u)=\sum_{s=1}^m\sum_{b=1}^s
               C_{d;s,b}(u)\clipbin{n-s+1}{b}.
\]
Every $C_{d;s,b}$ has nonnegative integer coefficients and is divisible
by $u^b$. Put $C_{s,b,k}=[u^k]C_{d;s,b}(u)$. Then
\[
 Q_{d,k}(x)=\sum_{j=0}^k A_{k,j}\binom{x-m+1}{j},\qquad
 A_{k,j}=\sum_{s=1}^m\sum_{b=\max(1,j)}^s
            C_{s,b,k}\binom{m-s}{b-j}\in\Z_{\ge0}.
\]
In particular, $\Delta^jQ_{d,k}(n)\ge0$ for every integer $n\ge m-1$
and every $j\ge0$. The empty-word case is separately $L_{d,n}=1$.
\end{theorem}

\begin{proof}
Let $p$ be a feasible outcome and list its distinct used sites as
$x_1<\cdots<x_s$. Replace every occurrence of $x_a$ by $a$, giving a
unique surjective proper word $v$. Set
\[
 C=\{a\in[s-1]:x_{a+1}\ge x_a+2\}.
\]
The cuts divide the labels into $b=|C|+1$ runs of consecutive used
sites. The active consecutive block immediately below $p_i=x_{v_i}$
has exactly $\beta_i(v,C)$ sites. Consequently its marked fiber is
$\prod_i(u+\beta_i(v,C))$, independently of the lengths of the unused
gaps.

Conversely, $v$, $C$, and a placement of their runs uniquely recover
the outcome. The total length of the $b$ runs is $s$. Write $g_0,g_b$
for the two exterior gaps and $g_1,\ldots,g_{b-1}$ for the interior
gaps. The first two are nonnegative; the interior gaps are at least
one; and their sum is $n-s$. Subtracting the mandatory $b-1$ unused
sites leaves $b+1$ nonnegative variables with sum $n-s-b+1$.
Stars and bars therefore gives $\binom{n-s+1}{b}$ placements when
$n\ge s+b-1$, and none otherwise. This proves the packed formula.

For each run, let $\ell$ be its least label. Any occurrence of this
label has $\beta_i=0$: either $\ell=1$, or $\ell-1$ is a cut.
Surjectivity supplies at least one such occurrence in each of the
$b$ different runs, so each product is divisible by $u^b$.

When $n\ge m-1$, every upper argument $n-s+1$ is a nonnegative
integer. The clipped binomial then equals its binomial polynomial,
including when the upper argument is smaller than $b$. The identity
\[
 \binom{x-s+1}{b}
   =\sum_{j=0}^b\binom{x-m+1}{j}\binom{m-s}{b-j}
\]
is Vandermonde's identity for binomial polynomials. Taking $[u^k]$
and using $C_{s,b,k}=0$ for $b>k$ gives the stated Newton coefficients.
The resulting polynomial agrees with the stable layer on every
sufficiently large integer and hence equals $Q_{d,k}$ identically.
Finally,
\[
 \Delta^j\binom{x-m+1}{\ell}
   =\begin{cases}\binom{x-m+1}{\ell-j},&\ell\ge j,\\0,&\ell<j,
     \end{cases}
\]
and these quantities are nonnegative at integers $x\ge m-1$.
\end{proof}

\begin{remark}The positive coefficients do not imply real-rootedness in $u$: $L_{4,3,1}(u)=u^2(7+24u+24u^2)$ has a quadratic factor with discriminant $-96$.\end{remark}

\subsection{An exact algorithm at fixed occupancy width}

Only departure indices through time $m$ affect the count, so first
replace $d_j$ by $\min(d_j,m)$. Put $B_0=c_0=1$ and
\[
 c_R=2^R\sum_{h=0}^R\frac{R!}{(R-h)!}.
\]

\begin{proposition}[Affine frontier algorithm]\label{prop:algorithm}
For $m\ge1$ and $n\ge1$, the full polynomial $L_{d,n}(u)$ can be
computed exactly with at most $B_R(2m+1)^{2R}$ states at each stored
time and at most $c_R$ transitions per state. Here $B_R$ is the $R$th
Bell number. An arithmetic cost bound is
\[
 O\bigl(mB_R(2m+1)^{2R}c_R(R^2+m+1)\bigr).
\]
The state space is independent of $n$, and bit complexity is
polynomial in $m$ and $\log(n+1)$ for fixed $R$. The same state space
also computes $\sum_kQ_{d,k}(x)u^k$ in polynomial time for fixed $R$.
For $m=0$ the output is one without any transitions.
\end{proposition}

\begin{proof}
Immediately before arrival $i$, the frontier is the known set $J_i$,
which has at most $R$ vertices. Process the local affine motifs in
chronological order, discarding inconsistent partial systems. A
stored state consists of the following data for each partial
constraint component meeting the frontier: its active vertices, their
integer potentials normalized to zero at the least active vertex,
and the minimum and maximum potentials over \emph{all} vertices of
the component, including forgotten vertices. Attach a polynomial in
$u$ whose coefficients are the signed sums of all histories with
these data; closed components have already contributed their
numerical translation factors. The initial empty state has weight
one.

To process arrival $i$, temporarily introduce its vertex as a
singleton with potential zero and choose one local motif. There are
$c_{r_i}\le c_R$ choices. An equation within a component is either
consistent with its potentials or discards the transition. An
equation between components fixes their relative translation;
translate all their recorded potentials and extrema and merge them.
Multiply the weight by the local sign and by $u$ for an empty
displacement decoration. Then forget all vertices whose truncated
departure index equals $i$. For a component with no remaining active
vertex, multiply by $(n-\sigma)_+$, where $\sigma$ is its recorded
span, and remove it. For each surviving component renormalize at its
least remaining vertex. After time $i<m$ the frontier is exactly
$J_{i+1}$. After time $m$, every component is closed, so the final
state is empty. In particular no state with the temporary $R+1$
vertices needs to be stored.

A closed component cannot meet any later equation: a departed vertex
never becomes active again. The retained potentials determine every
possible future consistency test and relative translation, and the
retained extrema determine its final translation count. Thus histories
with the same state have identical future continuations. The above
operations group, without changing their total, precisely the affine
terms of Theorem~\ref{thm:affine}.

The unit-chain argument applies also to every partial component, so
its span is at most $m-1$. Relative to an active anchor, all recorded
potentials and extrema belong to $[-m+1,m-1]$. If the frontier has
$r\le R$ vertices partitioned into $a$ components, the anchors remove
$a$ potential coordinates; there are $r-a$ other active potentials
and $2a$ extrema. There are consequently at most
\[
 (2m+1)^{r-a+2a}=(2m+1)^{r+a}\le(2m+1)^{2R}
\]
potential assignments for each partition. There are at most $B_R$
partitions, including the empty partition when $R=0$.

Component updates require $O(R^2)$ arithmetic operations per local
motif. Form the product of all numerical closing factors as one
scalar in $O(R+1)$ operations, then multiply the polynomial weight by
that scalar and perform its possible shift by $u$ in $O(m+1)$
operations. This proves the arithmetic bound. There are at most
$c_R^m$ indexed histories, each with a numerical weight of absolute
value at most $n^m$. Hence every intermediate integer coefficient
has bit length
\[
 O\bigl(m(\log(n+1)+\log(c_R+1))\bigr).
\]
The state integers have $O(\log(m+1))$ bits. Arithmetic and state
handling therefore have polynomial bit cost at fixed $R$.

For the stable polynomial, run the same transitions with each closing
factor replaced by $x-\sigma$, without pruning at numerical zeros.
Weights then have at most $(m+1)^2$ coefficients. The absolute sum
of the coefficients contributed by each history is at most $m^m$,
so their bit lengths are
$O(m(\log(m+1)+\log(c_R+1)))$. This is again polynomial for fixed
$R$, and the affine formula identifies the output with the stable
bivariate polynomial. This is an XP bound: its exponent may depend
on $R$.
\end{proof}


\section{Extreme lucky layers for common meters}


For $m\ge1$, write $N_{m,n,t,k}=[u^k]L_{m,n,t}(u)$. With $r_i=\min(t,i-1)$, $t\ge1$, direct counting gives
\[
N_{m,n,t,1}=\left(\prod_{i=2}^m r_i\right)(n-m+1)_+,\qquad
N_{m,n,t,m}=\prod_{i=1}^m(n-r_i)_+.
\]

For one lucky car, all outcomes are consecutive from $a_1$ and each later car chooses one of its $r_i$ active occupied preferences. For all lucky cars, each arrival has $n-r_i$ free choices. If exactly one car is unlucky, a division-free formula valid for every $n\ge1$ is
\[
N_{m,n,t,m-1}=\sum_{i=2}^m r_i\prod_{\substack{1\le j\le m\\j\ne r_i}}(n-r_j)_+.
\]

The omitted index is the numerical time index $r_i$, not $i$. To prove the formula when $n>\min(t,m-1)$, the all-lucky active $r_i$-tuple is uniform over injections. For a uniform occupied $r$-subset, its terminal suffix $L$ has mean $r/(n-r+1)$, by tail summation. The mean successful unlucky choices are $r-\E L=r(n-r)/(n-r+1)$. Multiply all-lucky prefix and suffix counts and cancel the factor $n-r_i+1=n-r_{r_i}$. In the remaining range every summand retains a zero positive-part factor. Finally $t\ge m-1$ gives $Q_{m,t}(x)=(x-m+1)(x+1)^{m-1}$.

\section{Compact enumeration when the meter is one}
\label{sec:one}

\subsection{A rational generating function and a positive recurrence}

Write $f_m^{(n)}=|\operatorname{MPF}_{m,n}(1)|$ for $n\geq1$ and $m\geq0$,
where the empty preference word contributes $f_0^{(n)}=1$.
The following formulas apply to arbitrarily many cars, including $m>n$.

\begin{theorem}[One-meter enumeration]\label{thm:one}
For every integer $n\geq1$,
\begin{equation}\label{eq:one-meter-gf}
 F_n(z):=\sum_{m\geq0}f_m^{(n)}z^m
 =\frac{(1+z)^n}{D_n(z)},\qquad
 D_n(z)=(1+z)^n(1-nz+z^2)-z^{n+2}.
\end{equation}
For $n\geq2$, the displayed fraction is reduced and its denominator has
degree exactly $n$. More explicitly,
\begin{equation}\label{eq:one-meter-positive-denominator}
 D_n(z)=1-\sum_{j=2}^{n}c_{n,j}z^j,
 \qquad
 c_{n,j}=\binom{n+2}{j}\frac{(j-1)(n+1-j)}{n+1}
 \in\mathbb Z_{>0}.
\end{equation}
Consequently,
\begin{equation}\label{eq:one-meter-positive-recurrence}
 f_m^{(n)}=\sum_{j=2}^{n}c_{n,j}f_{m-j}^{(n)}
 \qquad(m>n),
\end{equation}
with initial values
\begin{equation}\label{eq:one-meter-initial-values}
 f_m^{(n)}=
 \sum_{j=0}^{\lfloor m/2\rfloor}
 (-1)^j\binom{m-j}{j}n^{m-2j}
 \qquad(0\leq m\leq n).
\end{equation}
There is also a finite formula valid for every $m\geq0$:
\begin{equation}\label{eq:one-meter-multinomial}
 f_m^{(n)}=
 \sum_{\substack{0\leq b\leq n,\ k_2,\ldots,k_n\geq0\\
 b+\sum_{j=2}^{n}jk_j=m}}
 \binom nb
 \frac{(\sum_{j=2}^{n}k_j)!}{\prod_{j=2}^{n}k_j!}
 \prod_{j=2}^{n}c_{n,j}^{k_j}.
\end{equation}
When $n=1$, the tuples $(k_2,\ldots,k_n)$ are empty and all their
products have value one. In this case $D_1(z)=1$, $F_1(z)=1+z$, and
$f_m^{(1)}=0$ for every $m\geq2$; the empty sum in
\eqref{eq:one-meter-positive-recurrence} gives the same conclusion.
\end{theorem}

\begin{proof}
At meter one, only the preceding car is present when the next car arrives.
If that car occupies $x$, a new car parking at $y\neq x$ can prefer $y$.
It has precisely one additional preference, namely $x$, when $y=x+1$.
Thus the transition weight between outcomes is
\[
 A_{x,y}=\mathbf1_{\{y\neq x\}}+\mathbf1_{\{y=x+1\}}.
\]
Let $\boldsymbol1$ be the all-one column vector, let
$J=\boldsymbol1\boldsymbol1^{\mathsf T}$, and let $N$ be the $n\times n$
matrix with ones on the superdiagonal and zeros elsewhere. Then
$A=J-I+N$. The first car has one preference for each possible outcome, so
\[
 f_m^{(n)}=\boldsymbol1^{\mathsf T}A^{m-1}\boldsymbol1
 \quad(m\geq1),\qquad
 F_n(z)=1+z\boldsymbol1^{\mathsf T}(I-zA)^{-1}\boldsymbol1.
\]
These are identities of formal power series.

Put $C=(1+z)I-zN$. Since $N^n=0$,
\[
 C^{-1}=\sum_{a=0}^{n-1}\frac{z^aN^a}{(1+z)^{a+1}},\qquad
 s_n(z):=\boldsymbol1^{\mathsf T}C^{-1}\boldsymbol1
 =\sum_{a=0}^{n-1}\frac{(n-a)z^a}{(1+z)^{a+1}}.
\]
Writing $q=z/(1+z)$ and using
\[
 \sum_{a=0}^{n-1}(n-a)q^a
 =\frac{n-(n+1)q+q^{n+1}}{(1-q)^2}
\]
gives
\[
 s_n(z)=n-z+\frac{z^{n+1}}{(1+z)^n}.
\]
Now $I-zA=C-z\boldsymbol1\boldsymbol1^{\mathsf T}$. The rank-one
inverse identity, applicable since $1-zs_n(z)$ has constant term one, yields
\[
 \boldsymbol1^{\mathsf T}(I-zA)^{-1}\boldsymbol1
 =\frac{s_n(z)}{1-zs_n(z)}.
\]
Therefore $F_n(z)=1/(1-zs_n(z))$, proving \eqref{eq:one-meter-gf}.

Use the convention $\binom nr=0$ for $r\notin\{0,\ldots,n\}$.
Before subtracting $z^{n+2}$, the coefficient of $z^j$ in the denominator is
\[
 d_{n,j}=\binom nj-n\binom n{j-1}+\binom n{j-2}.
\]
The coefficient at $j=n+1$ is $-n+n=0$, and the coefficient at $j=n+2$
is one, which is removed by the final subtraction. Also $d_{n,0}=1$,
$d_{n,1}=0$, and, for $2\leq j\leq n$,
\[
 -d_{n,j}=\binom{n+2}{j}\frac{(j-1)(n+1-j)}{n+1}.
\]
This proves \eqref{eq:one-meter-positive-denominator}. Integrality follows
from the expression for $d_{n,j}$, and the last displayed factors are
strictly positive. For $n\geq2$, the coefficient of $z^n$ is
$1-\binom{n+1}{2}\neq0$, so $\deg D_n=n$.
Moreover $D_n(-1)=-(-1)^{n+2}\neq0$.
The only possible common factor with $(1+z)^n$ is $1+z$, so the fraction is
reduced.

Comparing coefficients in $D_n(z)F_n(z)=(1+z)^n$ above degree $n$
proves \eqref{eq:one-meter-positive-recurrence}.
To compute the initial coefficients, rewrite
\eqref{eq:one-meter-gf} as
\[
 F_n(z)=
 \frac1{1-nz+z^2-z^{n+2}/(1+z)^n}
 \equiv\frac1{1-nz+z^2}\pmod{z^{n+2}}.
\]
Expanding $\sum_{\ell\geq0}(nz-z^2)^\ell$, its coefficient of $z^m$ is
\[
 \sum_{j=0}^{\lfloor m/2\rfloor}
 (-1)^j\binom{m-j}{j}n^{m-2j},
\]
which proves \eqref{eq:one-meter-initial-values}. In fact this equality
also holds for $m=n+1$, although the recurrence only needs the stated
initial range. Finally expand
\[
 F_n(z)=(1+z)^n
 \sum_{K\geq0}\left(\sum_{j=2}^{n}c_{n,j}z^j\right)^K.
\]
Choose the numerator degree $b$ and collect the multiplicities $k_j$ of
the selected denominator terms. The multinomial theorem gives
\eqref{eq:one-meter-multinomial}. For fixed $m$ only finitely many tuples
satisfy its degree constraint. The $n=1$ case follows either from
\eqref{eq:one-meter-gf} or directly because a second car cannot park
before the first leaves.
\end{proof}

\begin{remark}
For $n=2$, the formulas reduce to
\[
 F_2(z)=\frac{(1+z)^2}{1-2z^2},\qquad
 f_0^{(2)}=1,\quad f_1^{(2)}=2,\quad f_2^{(2)}=3,
\]
and $f_m^{(2)}=2f_{m-2}^{(2)}$ for $m>2$.
The initial formula is associated with the short-range recurrence in \cite{DH}; the rational function
and positive recurrence extend the enumeration to every $m>n$.
\end{remark}

\subsection{Specified lucky sets and their exact polynomial onset}

Fix $m\geq1$, and let $N_K(n)$ count preference words having lucky set
exactly $K\subseteq[m]$ at meter one. If $1\notin K$, then $N_K(n)=0$
for every $n\geq1$. Otherwise write
\[
 K=\{v_1=1<v_2<\cdots<v_k\},\qquad
 v_{k+1}=m+1,\qquad s_j=v_{j+1}-v_j-1.
\]
Thus $s_j$ is the number of unlucky cars following the $j$th lucky car
before the next lucky arrival, and
$q:=\sum_{j=1}^ks_j=m-k$.
Let $\mathcal I_k$ be the partitions of $[k]$ into consecutive nonempty
intervals, and put $s(B)=\sum_{j\in B}s_j$.
For real $x$, write $(x)_+=\max\{x,0\}$.

\begin{theorem}[Lucky-run formula]\label{thm:lucky-run-formula}
For every integer $n\geq1$,
\begin{equation}\label{eq:lucky-run-formula}
 N_K(n)=\sum_{\pi\in\mathcal I_k}(-1)^{k-|\pi|}
              \prod_{B\in\pi}(n-s(B))_+.
\end{equation}
Equivalently, let $H_0=1$ and set
\begin{equation}\label{eq:lucky-run-dp}
 H_j=\sum_{a=1}^{j}(-1)^{j-a}H_{a-1}
        \left(n-\sum_{h=a}^{j}s_h\right)_+
        \quad(1\leq j\leq k).
\end{equation}
Then $N_K(n)=H_k$. With the partial sums of the $s_j$ precomputed,
this uses $O(k^2)$ arithmetic operations.
The polynomial
\begin{equation}\label{eq:lucky-run-polynomial}
 Q_K(x)=\sum_{\pi\in\mathcal I_k}(-1)^{k-|\pi|}
              \prod_{B\in\pi}(x-s(B))
\end{equation}
satisfies $N_K(n)=Q_K(n)$ for every integer
$n\geq\max\{1,q\}$. It is monic of degree $k$, and
\begin{equation}\label{eq:lucky-run-second-coefficient}
 [x^{k-1}]Q_K(x)=-(m-1).
\end{equation}
For $m=0$, the same formulas hold with $K=\varnothing$, $k=0$,
one empty interval partition, and empty products equal to one;
then $N_{\varnothing}(n)=Q_{\varnothing}(n)=H_0=1$.
The coefficient assertion \eqref{eq:lucky-run-second-coefficient}
is asserted only when $m\geq1$.
\end{theorem}

\begin{proof}
An unlucky car must prefer the preceding car's occupied spot and then park
one position to its right. If the $j$th lucky car parks at $b_j$, its run
therefore has successive outcomes
$b_j,b_j+1,\ldots,b_j+s_j$.
Its successful choices are precisely $1\leq b_j\leq n-s_j$.
At the next lucky arrival, the only occupied spot is $b_j+s_j$;
the additional restriction is consequently
$b_{j+1}\neq b_j+s_j$.
These data determine the preference word uniquely: each lucky car prefers
$b_j$, and each unlucky car prefers the preceding outcome.

Apply inclusion--exclusion to these $k-1$ forbidden equalities.
A selected equality joins two neighboring run indices. Selected joins
therefore determine a partition $\pi\in\mathcal I_k$ and number
$k-|\pi|$. For a block $B=\{a,\ldots,b\}$, the equalities give
\[
 b_h=b_a+\sum_{j=a}^{h-1}s_j\quad(a\leq h\leq b).
\]
All lower and upper bounds on the $b_h$ reduce to
$1\leq b_a\leq n-\sum_{j=a}^{b}s_j$ because every $s_j$ is
nonnegative. Thus the block contributes $(n-s(B))_+$, independently of
the other blocks. This proves \eqref{eq:lucky-run-formula}.
Separating each partition of $[j]$ by its last block $\{a,\ldots,j\}$
gives \eqref{eq:lucky-run-dp}. The count of summands is
$\sum_{j=1}^kj=k(k+1)/2$.

Every block weight is at most $q$, so for $n\geq\max\{1,q\}$ the
positive parts in \eqref{eq:lucky-run-formula} can be removed.
The singleton partition is the only term of degree $k$, and contributes
leading coefficient one and next coefficient $-q$.
There are exactly $k-1$ partitions with $k-1$ blocks, each having one
join and contributing $-1$ to the coefficient of $x^{k-1}$.
All remaining partitions have degree at most $k-2$.
Therefore that coefficient is $-q-(k-1)=-(m-1)$.
For $k=1$, the same argument reads $Q_K(x)=x-q$, giving its constant
coefficient $-q=-(m-1)$ directly.
\end{proof}

\begin{corollary}[Exact onset]\label{cor:one-meter-exact-onset}
For $m\geq1$ and $1\in K$, let $N_*(K)$ be the smallest positive integer
$N$ such that $N_K(n)=Q_K(n)$ for every integer $n\geq N$.
If $q>0$, let $a$ and $b$ be the first and last indices with $s_j>0$.
Then
\begin{equation}\label{eq:one-meter-exact-onset}
 N_*(K)=
 \begin{cases}
 q,&q\geq3,\\
 2,&q=2\text{ and }s_1s_k>0,\\
 1,&q\in\{0,1\},\text{ or }q=2\text{ and }s_1s_k=0.
 \end{cases}
\end{equation}
For $q\geq2$, the discrepancy immediately below the general span bound is
\begin{equation}\label{eq:one-meter-onset-discrepancy}
 N_K(q-1)-Q_K(q-1)
 =(-1)^{b-a}(q-2)^{a-1+k-b},
\end{equation}
where $0^0=1$.
If $1\notin K$, or if $m=0$ and $K=\varnothing$, the onset is one.
\end{corollary}

\begin{proof}
Equality holds at every $n\geq\max\{1,q\}$ by the preceding theorem.
This already proves onset one when $q=0$ or $q=1$.
Suppose $q\geq2$ and evaluate at $n=q-1\geq1$.
A block factor differs with and without its positive part precisely when
its weight equals $q$. Such a block contains all positive run weights,
hence contains $\{a,\ldots,b\}$, and is unique in its partition.
Its ordinary factor is $-1$ and its truncated factor is zero.
Every other block has weight zero and contributes $n$.

The $b-a$ joins inside $\{a,\ldots,b\}$ are forced. Of the
$v:=a-1+k-b$ remaining joins, select $d$; there are $\binom vd$ choices
and $v-d$ remaining blocks outside the weight-$q$ block.
The exact count minus the polynomial is thus
\[
 (-1)^{b-a}\sum_{d=0}^{v}\binom vd(-1)^d n^{v-d}
 =(-1)^{b-a}(n-1)^v.
\]
Substituting $n=q-1$ gives \eqref{eq:one-meter-onset-discrepancy}.
For $q\geq3$ this is nonzero, so onset $q$ is necessary and sufficient.
For $q=2$, it is nonzero exactly when $v=0$, equivalently $a=1$ and
$b=k$, or $s_1s_k>0$. In that case the onset is two.
Otherwise equality also holds at $n=1$, and hence at every positive
integer $n$, so the onset is one.
The identically zero and empty-word cases are immediate.
\end{proof}

\section{Boundary failure and overlap graphs}\label{sec:boundary}

Let $S$ denote success through $m$ arrivals, and set
\[
 C=\{\text{no overlap edge }\{j,i\}\text{ has }a_j=a_i=n\},
\]
\[
 E=\sum_i r_i=|E(G_d)|,\qquad
 \eta=\sum_v\deg_G(v)(\deg_G(v)-1),
\]
\[
 W=\sum_i\left(\sum_{j\in J_i}r_j+\binom{r_i}{2}\right),
 \qquad h=\min(n,R+2).
\]

\begin{theorem}[Finite graph comparison]\label{thm:graph}
For independent uniform preferences,
\[
 S\subseteq C,\qquad 0\le\Prob(C)-\Prob(S)\le W(h/n)^3.
\]
For $n\ge2$, writing $I_G(z)=\sum_{A\text{ independent}}z^{|A|}$,
\[
 \Prob(C)=(1-1/n)^mI_G\bigl(1/(n-1)\bigr).
\]
For $E>0$,
\[
 \left|\Prob(C)-e^{-E/n^2}\right|
 \le\min(1,n^2/E)
       \left(\frac{E+\eta}{n^4}+\frac{\eta}{n^3}\right).
\]
\end{theorem}

\begin{proof}
A successful car with preference $n$ occupies $n$ until its scheduled
departure. Two overlapping cars with this preference cannot both
succeed, proving $S\subseteq C$.

On $C\setminus S$, let $i$ be the first failed arrival. Every earlier
successful car satisfies $0\le p_j-a_j\le r_j\le R$. If $a_i=n$,
let $j\in J_i$ occupy $n$. Since the direct pair is absent,
$a_j<n$. When $j$ parked, some $k\in J_j$ occupied $n-1$; these are
three distinct cars $k<j<i$. There are $\sum_i\sum_{j\in J_i}r_j$
possible indexed witnesses. If $a_i<n$, both $n-1$ and $n$ are
occupied by two distinct cars in $J_i$, giving at most
$\sum_i\binom{r_i}{2}$ witnesses of the other type. In either case
all three preferences belong to
\[
 T_n=\{\max(1,n-R-1),\ldots,n\},\qquad |T_n|=h.
\]
Indeed a car occupying $n-1$ or $n$ has preference at least
$n-R-1$, while a failed scan has preference at least $n-R$.
The indexed cars have independent preferences, so each witness
event has probability $(h/n)^3$. A union bound proves the comparison.
When $n=1$, every possible failure is already a direct pair, and the
same bound holds.

The set $\{i:a_i=n\}$ is a Bernoulli-$1/n$ vertex set. It satisfies
$C$ precisely when it is independent, giving the independent-set
identity. For each edge put $Z_e=\ind\{a_i=a_j=n\}$ and let its
dependency neighborhood contain itself and all edges sharing an
endpoint. There are $\eta$ ordered pairs of distinct neighboring
edges. Thus the Chen--Stein quantities are
\[
 b_1=(E+\eta)n^{-4},\qquad b_2=\eta n^{-3},\qquad b_3=0,
 \qquad\lambda=E/n^2.
\]
The zero-event bound in \cite{AGG} proves the final inequality.
\end{proof}

\begin{corollary}[Highest stable coefficients]\label{cor:top}
For every schedule of $m\ge2$ cars,
\[
 [x^m]Q_d=1,\qquad [x^{m-1}]Q_d=0,
 \qquad [x^{m-2}]Q_d=-E.
\]
\end{corollary}
\begin{proof}
Fix the schedule and let $n\to\infty$. Each direct terminal-edge
event has probability $n^{-2}$. Two distinct such events jointly
fix at least three terminal preferences, so inclusion--exclusion
gives $\Prob(C)=1-E/n^2+O(n^{-3})$. The witness comparison also has
error $O(n^{-3})$. Consequently
\[
 L_{d,n}(1)=n^m-En^{m-2}+O(n^{m-3}).
\]
Eventual polynomiality identifies the three coefficients. For $m=2$
one can alternatively read the exact counts $n^2$ and $n^2-1$.
\end{proof}

\begin{corollary}[The first coefficient detecting chronology]\label{cor:third}
For a schedule of $m\ge3$ cars, put
\[
 o_j=|\{i>j:j\in J_i\}|,\qquad
 A=\sum_j\binom{o_j}{2},\qquad T=\sum_i\binom{r_i}{2}.
\]
Then
\[
 [x^{m-3}]Q_d(x)=A-3T.
\]
For a common meter, $o_j=\min(t,m-j)$ and hence $A=T$; its coefficient
is $-2\sum_i\binom{\min(t,i-1)}2$.
\end{corollary}

\begin{proof}
For $a<b<c$, an overlap edge $\{a,c\}$ forces $\{a,b\}$, because
$d_a\ge c>b$. The possible connected induced three-vertex graphs
are therefore an old-centered fork (edges $ab,ac$), a chronological
chain (edges $ab,bc$), and a triangle. Denote their numbers by $F$,
$P$, and $T$, respectively. The definition of $T$ indeed counts
triangles, since every $J_i$ is a clique. Inclusion--exclusion for
terminal-edge events gives, at fixed schedule,
\[
 \Prob(C)=1-\frac E{n^2}
                 +\frac{F+P+2T}{n^3}+O_m(n^{-4}).
\]
Each two-edge induced triple contributes one pair intersection. A
triangle contributes three pair intersections and one triple
intersection, for net contribution two. All intersections involving
at least four vertices contribute $O_m(n^{-4})$.

We next determine failure outside $C$ to the same accuracy. At the
first failed arrival $i$, recursively collect all cars occupying a
site scanned by $i$, and, for each collected earlier car $j$, all
cars occupying its scanned sites $a_j,\ldots,p_j-1$ when it parked.
Include the failed car $i$ itself and call the resulting set its causal closure. Every ancestor has
strictly smaller time index. Retain the original time labels and
departure indices when deleting cars. By chronological induction,
deleting all cars outside the closure preserves every successful outcome in
the closure and preserves its final failed scan: all occupied sites
required by each scan were included, while its successful target
was free originally and remains free after deletion.

A failed scan satisfies $a_i\ge n-R$. Along a backward dependency
step to a blocker $j$, its outcome is at least the child's
preference and its displacement is at most $R$. Hence the preference
decreases by at most $R$ per step. A chain has at most $m-1$ steps,
so every closure preference is at least $n-mR\ge n-m(m-1)$.
For sufficiently large $n$, all of them therefore lie in the fixed
terminal window of $M=m^2+1$ sites. The probability that four or more
cars have preferences in this window is $O_m(n^{-4})$.

Outside this exceptional event, a failure on $C$ has a causal closure
of exactly three cars. One car cannot fail, and a closure of two
cars can fail only from an overlapping pair with both preferences
$n$, excluded by $C$. Its three-vertex graph is connected and is
therefore one of the three patterns above. A direct three-car check
gives the failures on $C$: a chronological chain has precisely
$(a_a,a_b,a_c)=(n-1,n-1,n)$; an old-centered fork has none; a triangle
has the all-$(n-1)$ assignment and the three assignments containing
one $n$ and two $n-1$'s. These counts are $1,0,4$.

Conversely, each listed assignment causes a failure in the full
process if every other preference is outside the $M$-site window.
Such a car has outcome at most $n-M+R<n-1$ and cannot affect any
scan on $\{n-1,n\}$. Thus the triple behaves exactly as it does in
the restricted process. Its probability is $n^{-3}+O_m(n^{-4})$.
Different triples can overlap only on an event with at least four
window preferences; different assignments on the same triple are
disjoint. We conclude that
\[
 \Prob(C)-\Prob(S)=\frac{P+4T}{n^3}+O_m(n^{-4}),
 \qquad
 \Prob(S)=1-\frac E{n^2}+\frac{F-2T}{n^3}+O_m(n^{-4}).
\]
Every choice of two later neighbors of a vertex is either an
old-centered fork or a triangle, so $A=F+T$. Multiplying by $n^m$
and using eventual polynomiality yields $A-3T$. Finally the
common-meter in-degree and out-degree sequences are reversals of
one another, proving $A=T$.
\end{proof}


\begin{example}
For $d=(2,3,3)$, the chronological path has $E=2$ and $A=T=0$, giving $Q_d(x)=x^3-2x$. For $d=(3,2,3)$, the oldest-centered star has $E=2$, $A=1$, and $T=0$, giving $Q_d(x)=x^3-2x+1$. The two abstract overlap graphs are isomorphic, but their chronological roles give different constant coefficients.
\end{example}

\clearpage
\section{Further questions}
For common meters the component-span bound is attained, but cancellation in the affine sum can make the actual polynomial onset smaller. Meter one admits a complete onset classification through its interval-partition formula. Determining the corresponding onset for larger meters remains open. It would also be useful to find compact closed formulas for other fixed meters, and to decide whether the exact frontier computation can be improved from an XP bound to a running time $f(R)m^C$ with an exponent $C$ independent of the occupancy width.

The overlap graph determines the leading lucky-layer data and the first boundary coefficients; Corollary~\ref{cor:third} shows how the next coefficient also uses the chronological in-degrees and out-degrees. It does not determine the full enumeration: different chronological realizations of the same abstract graph can give different counts. A direct combinatorial description of the lower stable coefficients would clarify how these two kinds of information interact.

\begingroup
\small
\setlength{\parskip}{0pt}
\begin{thebibliography}{99}
\bibitem{DH} S. Daugherty, P. E. Harris, I. Klein, and M. McClinton, \emph{Metered Parking Functions}, Integers \textbf{25} (2025), A73. DOI: 10.5281/zenodo.16881806; arXiv:2406.12941.
\bibitem{FZ} D. Forge and T. Zaslavsky, \emph{Lattice point counts for the Shi arrangement and other affinographic hyperplane arrangements}, J. Combin. Theory Ser. A \textbf{114} (2007), 97--109. DOI: 10.1016/j.jcta.2006.03.006.
\bibitem{BC} P. Berthom\'e, R. Cordovil, D. Forge, V. Ventos, and T. Zaslavsky, \emph{An elementary chromatic reduction for gain graphs and special hyperplane arrangements}, Electron. J. Combin. \textbf{16} (2009), R121. DOI: 10.37236/210.
\bibitem{HM} P. E. Harris and L. Martinez, \emph{Parking functions with a fixed set of lucky cars}, arXiv:2410.08057.
\bibitem{FHMS} M. Ferreri, P. E. Harris, L. Martinez, and E. Swartz, \emph{Enumerating vector parking functions and their outcomes based on specified lucky cars}, arXiv:2508.13917.
\bibitem{AGG} R. Arratia, L. Goldstein, and L. Gordon, \emph{Two Moments Suffice for Poisson Approximations: The Chen--Stein Method}, Ann. Probab. \textbf{17} (1989), 9--25. DOI: 10.1214/aop/1176991491.
\end{thebibliography}
\endgroup
\end{document}
